%% kapa: exam=1 round=29 year=2018 period=1 form=A subject=real_estate subject_ko=부동산학원론 seq=11 number=91
%% source: Q-Net 감정평가사 기출문제 (cst00309) artlSeq=5202808 file=extracted/1008/5202808/2018년 제29회 감정평가사 1차시험 기출문제/2018년 제29회 감정평가사 1차시험 기출문제/1교시 A형.pdf
%% conversion: kordoc 4.24.51 -> tools/gen1.py (auto); review=proofread-sonnet; flags=tikz
\begin{question}{91}
D도시 인근에 A, B, C 세 개의 쇼핑센터가 있다. 허프(Huff)의 상권분석모형을 적용할 경우, 각 쇼핑센터의 이용객 수는? (단, 거리마찰계수: 2, D도시 인구의 40\%가 위 쇼핑센터의 이용객이고, A, B, C 중 한 곳에서만 쇼핑함)
\begin{center}
\begin{tikzpicture}[x=0.13mm,y=-0.13mm,font=\scriptsize,line width=0.6pt]
\draw (465,72) ellipse (118 and 70);
\node[align=center] at (465,72) {B 쇼핑센터\\매장면적: 10,000 ㎡};
\draw (467,142) -- (467,220);
\node at (432,181) {5Km};
\draw (500,290) ellipse (135 and 70);
\node[align=center] at (500,290) {D 도시\\인구: 60만명};
\draw (115,285) ellipse (118 and 70);
\node[align=center] at (115,285) {A 쇼핑센터\\매장면적: 4,000 ㎡};
\draw (233,287) -- (365,290);
\node at (300,262) {10Km};
\draw (880,150) ellipse (118 and 75);
\node[align=center] at (880,150) {C 쇼핑센터\\매장면적: 20,000 ㎡};
\draw (620,262) -- (790,198);
\node at (700,214) {10Km};
\end{tikzpicture}
\end{center}
\begin{choices}
  \item A: 15,000명, B: 150,000명, C: 75,000명
  \item A: 15,000명, B: 155,000명, C: 70,000명
  \item A: 15,000명, B: 160,000명, C: 65,000명
  \item A: 16,000명, B: 150,000명, C: 74,000명
  \item A: 16,000명, B: 155,000명, C: 69,000명
\end{choices}
\end{question}
\answer{1}
